\begin{table}[H]\centering\small
\caption{1月25--31日增益边界诊断；首末库存均为6000 kWh}
\label{tab:revision-gain}
\begin{tabular}{lrrrrr}\toprule
情境 & 增益列数 & 结构零列 & 有效列触边 & 界限10费用 & 界限20费用\\\midrule
问题二 & 768 & 388 & 0/380 & 347089.63 & 347089.63 \\
问题三 & 768 & 160 & 0/608 & 341917.12 & 341902.65 \\
问题四：日前 & 768 & 388 & 0/380 & 386128.80 & 386128.80 \\
问题四：日内 & 768 & 160 & 0/608 & 380647.04 & 380647.04 \\
\bottomrule\end{tabular}\end{table}

补充诊断在1月25--31日逐列检查约束矩阵与目标中的精确零列，将其从有效反馈分母剔除。两种界限下有效列触边均为零；问题三费用相差14.47元，仅0.00423\%，另三问不变。放宽界限可能影响非唯一最优解及浮点择优，不据此宣称策略类扩张产生了稳定收益，故不修改冻结的界限10。此表是开发期诊断，不能将其有效列触边比例外推为全年比例。
